\documentclass[10pt,a4paper]{amsart}
\usepackage[margin=19mm]{geometry}
\usepackage{amsmath,amssymb,mathtools}
\usepackage[scr=rsfs,cal=euler]{mathalfa}
\usepackage{xfrac}
\usepackage[unicode,hidelinks,bookmarks=false]{hyperref}
\newcommand{\ie}{i.e.\ }
\newcommand{\cf}{cf.\ }
\newcommand{\resp}{respectively\ }
\newcommand{\Subsection}{\subsection}
\providecommand{\nobreakdash}{\nobreak\hbox{-}}
\setlength{\emergencystretch}{2em}
\multlinegap0pt
\usepackage{amssymb}
\usepackage{tikz}
\newtheorem{thm}{Theorem}
\newtheorem{exa}{Example}
\newtheorem{lem}{Lemma}
\newcommand{\E}{\mathbf{E}}
\newcommand{\bR}{\mathbb{R}}
\newcommand{\ba}{\begin{equation}\begin{aligned}}
\newcommand{\di}{\mathrm{d}}
\newcommand{\e}{\varepsilon}
\newcommand{\ea}{\end{aligned}\end{equation}}
\title{Diffusion limits of cyclic finite-velocity random motions\\ along vector fields}
\author{Olga Aryasova, Ilya Pavlyukevich, Yusif Sualah}
\date{Source: arXiv:2608.22514v2 (CC BY 4.0)}
\begin{document}
\maketitle

\noindent\emph{Source and licence.} Diffusion limits of cyclic finite-velocity random motions\\ along vector fields. Version arXiv:2608.22514v2 (CC BY 4.0), distributed by arXiv under the Creative Commons Attribution 4.0 International licence, http://creativecommons.org/licenses/by/4.0/, as recorded in the arXiv licence metadata for this exact version and shown on the abstract page https://arxiv.org/abs/2608.22514v2. Full licence notice: \texttt{licenses/arxiv-2608.22514-CC-BY-4.0.txt}.\par\smallskip

\noindent\textit{Editorial scope.} Counted unit: the article's lem l:A (source0.tex lines 1346--1356). The article's complete original proof of it is delivered with it (source0.tex lines 1358--1633, one \texttt{proof} environment of 276 lines). No mathematics is written, repaired or re-derived: the statement and the proof are the article's own lines, in the article's own notation, and no retained line is substituted or re-flowed.  Retained uncounted as the source-local context the proof needs: lines 597--610; lines 614--646; lines 695--741; lines 743--749.  One declared adaptation: exactly 1 ancillary strings inside the retained spans are omitted (image-inclusion commands and, where the source repeats a label, the repeated label definitions) and no other line differs from the source; the figure environments, with their placement, captions and labels, are kept, and the package records the omitted strings in fidelity receipts bound to the affected bodies.  No retained line contains a citation, so the wrapper carries no bibliography and no bibliography entry.  The retained lines contain the article's own diagram code, which the wrapper typesets with the standard tikz packages; no image file is delivered and the package declares no asset dependency.  Editorial formatting only, in addition: an AMS \texttt{amsart} preamble that restates the article's own declarations and macros used by the retained lines, namely the environments \texttt{thm}, \texttt{exa}, \texttt{lem} on separate counters, together with the standard packages those lines need.  The retained environments print in this document as Theorem 1, Exa 1, Lemma 1 in order of appearance, whereas the article prints the same items with the numbering of its own full document.  The complete native source of the article as distributed in the arXiv e-print for this exact version is bundled unchanged as \texttt{sources/208339/source0.tex}.
For the formulation of the main result we will need the directional derivatives $D V_i[V_j]=DV_i(x)V_j(x)$ of the vector fields $V_i$
in direction $V_j$ defined as
vector fields with coordinates
\ba
\label{e:DVV}
(D V_i[V_j])^\alpha = \sum_{\beta=1}^d V_j^\beta   \cdot  \partial_\beta V_i^\alpha,\quad \alpha=1,\dots, d,\  1\leq i,j\leq p,
\ea
and the Lie bracket $[V_i,V_j]$ defined as
\ba
{} [V_i,V_j]&= D V_j [ V_i] - D V_i [ V_j],\quad 1\leq i,j\leq p,\\
([V_i,V_j])^\alpha &
= \sum_{\beta=1}^d  \Big(V_i^\beta  \cdot  \partial_\beta V_j^\alpha
- V_j^\beta \cdot \partial_\beta V_i^\alpha\Big) ,\quad \alpha=1,\dots, d,\  1\leq i,j\leq p.
\ea

\begin{thm}
\label{t:A}
Suppose that assumptions
$\mathbf{A}_V$,
$\mathbf{A}_\tau$,
$\mathbf{A}_{V,\mu}$ and
$\mathbf{A}_\varkappa$ hold. Then, for any initial values $x\in\bR^d$ and $\iota\in \{1,\dots,p\}$,
the flight evolutions $\{X^{\e,\delta}\}$ converge weakly in $C(\bR_+,\bR^d)$
to a diffusion $X$ starting at $x$ with the generator
\ba
\label{e:opA}
Af
= \frac{\varkappa^2}{2\mu}  \sum_{i} \sigma_{i}^2 \langle \nabla^2 f  V_i, V_i\rangle
+ \frac{\varkappa^2}{\mu} \sum_{j< i} \mu_i\mu_j \langle \nabla f , D V_i [ V_j] \rangle,\quad f\in C^3_b(\bR^d,\bR).
\ea
The diffusion $X$ is a solution to an It\^o-SDE
\ba
\label{e:SDEX}
\di X_t= \frac{\varkappa}{\sqrt\mu}  \sum_{i} \sigma_{i} V_i(X_t)\,\di W^i_t +
\frac{\varkappa^2}{\mu} \sum_{j< i} \mu_i\mu_j  D V_i [ V_j](X_t)\,\di t,
\ea
where $W^1,\dots,W^p$ are independent standard one-dimensional Brownian motions.
The SDE \eqref{e:SDEX} can be rewritten in the Stratonovich form
as
\ba
\label{e:XStrat}
\di X_t= \frac{\varkappa}{\sqrt\mu}  \sum_{i} \sigma_{i} V_i(X_t)\circ\di W^i_t
-\frac{\varkappa^2}{2\mu} \Big(
\sum_{i} \lambda_i^2  D V_i [ V_i](X_t)
+\sum_{j< i} \mu_i\mu_j  [ V_i,V_j](X_t)
\Big)\,\di t.
\ea
\end{thm}

\begin{exa}[cyclic motion along two vector fields in $\bR^d$]
\label{exa:UV}
Let $(T_k)_{k\in\mathbb N_0}$ be the arrival times of a Poisson process with intensity $c\in (0,\infty)$.
Then, $ T_{k+1}-T_k\stackrel{\di}{=} \mathrm{Exp}(c)$, $\mu_i=1/c$, $\lambda_i^2=2/c^2$,
$\sigma_i^2=1/c^2$, $i=1,2$.
Let $V,U\colon\bR^d\to\bR^d$ be two vector fields. Then, any permutation of the vector fields
$\{U,V,-U,-V\}$ defines a four-step cyclic motion in $\bR^d$, $p=4$.
Since a cyclic shift corresponds only to a change of the initial direction
and does not affect the limiting diffusion, the $4!=24$ permutations reduce
to six inequivalent cyclic orderings:
\ba
(V^{(1)}_1,V^{(1)}_2,V^{(1)}_3,V^{(1)}_4)&=(U, -U, V, -V),\\
(V^{(2)}_1,V^{(2)}_2,V^{(2)}_3,V^{(2)}_4)&=(U, -U, -V, V),\\
(V^{(3)}_1,V^{(3)}_2,V^{(3)}_3,V^{(3)}_4)&=(U, V, -U, -V),\\
(V^{(4)}_1,V^{(4)}_2,V^{(4)}_3,V^{(4)}_4)&=(U, V, -V, -U),\\
(V^{(5)}_1,V^{(5)}_2,V^{(5)}_3,V^{(5)}_4)&=(U, -V, -U, V),\\
(V^{(6)}_1,V^{(6)}_2,V^{(6)}_3,V^{(6)}_4)&=(U, -V, V, -U).
\ea
These orders are schematically represented in Fig.\ \ref{f:two}.
In particular, the sequences $(3)$ and $(5)$ encode the counterclockwise and clockwise motions, respectively.
The limiting operators $A$ and $B$ have the following form:
\ba
A^{(1)}f&=A^{(2)}f=A^{(4)}f =A^{(6)}f\\
&= \frac{\varkappa^2}{4c} \Big( \langle \nabla^2 f  V, V\rangle +\langle \nabla^2 f  U, U\rangle
- \langle \nabla f , D U [ U] + D V [ V]\rangle\Big),\\
A^{(3)}f&=\frac{\varkappa^2}{4c} \Big(
\langle \nabla^2 f  V, V\rangle +\langle \nabla^2 f  U, U\rangle
- \langle \nabla f , D U [ U] + D V [ V] -[U,V] \rangle\Big),\\
A^{(5)}f&=\frac{\varkappa^2}{4c} \Big(
\langle \nabla^2 f  V, V\rangle +\langle \nabla^2 f  U, U\rangle
- \langle \nabla f , D U [ U] + D V [ V] +[U,V] \rangle\Big),\\
\ea
and
\ba
B^{(1)}f&=B^{(2)}f=B^{(4)}f =B^{(6)}f\\
&= \frac{\varkappa^2}{4c} \Big( \langle \nabla^2 f  V, V\rangle +\langle \nabla^2 f  U, U\rangle
+\langle \nabla f , D U [ U] + D V [ V]\rangle\Big),\\
B^{(3)}f&=\frac{\varkappa^2}{4c} \Big( \langle \nabla^2 f  V, V\rangle +\langle \nabla^2 f  U, U\rangle
+\langle \nabla f , D U [ U] + D V [ V] +[U,V]\rangle\Big),\\
B^{(5)}f
&= \frac{\varkappa^2}{4c} \Big( \langle \nabla^2 f  V, V\rangle +\langle \nabla^2 f  U, U\rangle
+ \langle \nabla f , D U [ U] + D V [ V]-[U,V]\rangle
\Big).
\ea
In particular, the orderings (3) and (5), corresponding to opposite orientations of the cycle,
produce Lie-bracket drift terms with opposite signs.
Thus, the orientation of the microscopic cyclic motion remains visible in the macroscopic diffusion limit.

\begin{figure}
\begin{center}

\end{center}
\caption{Six inequivalent cyclic motions along vector fields $\pm U$ and $\pm V$.\label{f:two}}
\end{figure}
\end{exa}

\begin{lem}
\label{l:A}
For any $f\in C^3_b(\bR^d,\bR)$,
$x\in\bR^d$, any $\iota\in \{1,\dots,p\}$ we have
\ba
\label{e:Aiota}
\lim_{\e,\delta \downarrow 0} \sup_{x\in\mathbb R^d,\iota\in\{1,\dots,p\}}
|A_\iota^{\e,\delta }f(x) -A f(x)|=0,
\ea
where the operator $A$ is defined in \eqref{e:opA}.
\end{lem}

\begin{proof}
In order prove \eqref{e:Aiota} we first expand the contribution of the first-order Taylor terms of
$f(Z^{\e,\delta}_{T^\e_{p}})$
and use the centering condition to identify the order-dependent drift.
We then treat the second-order Taylor terms, which determine the diffusion coefficient.
Finally, we show that the resulting operator is invariant under cyclic shifts of the initial direction.

1. Let $\iota=1$.
For $i=1,\dots,p$,
denote $\Delta_i T^\e:=T_{i}^\e-T_{i-1}^\e\stackrel{\di}{=}\e \tau_i$. Then, the Taylor formula yields
\ba
\label{e:f}
\E_{x,1}f(Z^{\e,\delta}_{T_p^\e})& -f(x)
=\sum_{i=1}^{p} \E_{x,1}\Big[f(Z^{\e,\delta}_{T_{i}^\e})- f(Z^{\e,\delta}_{T_{i-1}^\e})\Big]\\
&= \sum_{i=1}^{p}\E_{x,1}
\Big[f \Big(Z^{\e,\delta}_{T_{i-1}^\e} + \frac{1}{\delta}V_{i}(Z^{\e,\delta}_{T_{i-1}^\e} )\Delta_i T^\e\Big)
-f (Z^{\e,\delta}_{T_{i-1}^\e} )\Big]\\
&= \frac{1}{\delta}\sum_{i=1}^{p}
\E_{x,1}
\langle \nabla f(Z^{\e,\delta}_{T_{i-1}^\e} ) , V_{i} (Z^{\e,\delta}_{T_{i-1}^\e} )\Delta_i T^\e\rangle\\
&+ \frac{1}{\delta^2} \sum_{i=1}^{p}
\E_{x,1} \int_0^1  \Big\langle \nabla^2 f\Big(Z^{\e,\delta}_{T_{i-1}^\e}
+  \frac{\theta}{\delta}V_{i} (Z^{\e,\delta}_{T_{i-1}^\e} )\Delta_i T^\e\Big) V_i (Z^{\e,\delta}_{T_{i-1}^\e} )\Delta_i T^\e,
V_{i} (Z^{\e,\delta}_{T_{i-1}^\e} )\Delta_i T^\e\Big\rangle(1-\theta)\,\di \theta.
\ea
a) We expand and estimate the first term in the last sum in \eqref{e:f}.
To exploit the centering condition in the first-order Taylor contribution, set
\ba
\label{e:g}
g_i(x):=\langle \nabla f(x),  V_i (x)\rangle.
\ea
Then $\mathbf{A}_{V,\mu}$ implies
\ba
\sum_{i=1}^{p } \mu_i g_i(x) = \Big\langle \nabla f(x), \sum_{i=1}^{p }\mu_i V_i(x)\Big\rangle \equiv 0,\quad
x\in\bR^d.
\ea
Consequently
we get
\ba
\frac{1}{\delta}\sum_{i=1}^{p} \E_{x,1} \Big[ g_i (Z^{\e,\delta}_{T_{i-1}^\e} ) \Delta_i T^\e\Big]
&=  \frac{\e}{\delta} \sum_{i=1}^{p} \mu_i \E_{x,1}\Big[ g_i(Z^{\e,\delta}_{T_{i-1}^\e} ) - g_i(x) \Big]\\
&=  \frac{\e}{\delta} \sum_{i=2}^{p}  \mu_i \E_{x,1}\Big[ g_i(Z^{\e,\delta}_{T_{i-1}^\e} ) - g_i(x) \Big]\\
&=   \frac{\e}{\delta}\sum_{i=2}^{p}\mu_i\sum_{j=1}^{i-1}   \E_{x,1}
\Big[ g_i(Z^{\e,\delta}_{T_{j}^\e} ) - g_i(Z^{\e,\delta}_{T^\e_{j-1}}) \Big].
\ea
For each $j<i$, we have
\ba
 \frac{\e}{\delta}\E_{x,1}[g_i(Z^{\e,\delta}_{T_{j}^\e} ) &- g_i(Z^{\e,\delta}_{T^\e_{j-1}})]
= \frac{\e}{\delta^2}\E_{x,1}\Big[ \Delta_j T^\e
\int_0^1 \langle \nabla
g_i(Z^{\e,\delta}_{T_{j-1}^\e}
+\frac{\theta \Delta_j T^\e}{\delta }
V_j(Z^{\e,\delta}_{T_{j-1}^\e} )), V_j(Z^{\e,\delta}_{T_{j-1}^\e} )\rangle\,\di \theta \Big]\\
&=  \frac{\e^2}{\delta^2} \mu_j \langle \nabla g_i(x), V_j(x) \rangle\\
&+
\frac{\e}{\delta^2}
\E_{x,1}\Big[\Delta_j T^\e
\int_0^1 \Big(\langle \nabla
g_i(Z^{\e,\delta}_{T_{j-1}^\e}
+ \frac{\theta \Delta_j T^\e}{\delta }V_j(Z^{\e,\delta}_{T_{j-1}^\e} )), V_j(Z^{\e,\delta}_{T_{j-1}^\e} ) \rangle
- \langle \nabla g_i(x), V_j(x)\rangle\Big)\,\di \theta \Big]\\
&= \frac{\e^2}{\delta^2} \mu_j \langle \nabla g_i(x), V_j(x) \rangle + \E_{x,1}R^{\e,\delta}_{ij}.
\ea
Overall,
\ba
\frac{1}{\delta}\sum_{i=1}^{p}
\E_{x,1}
\langle \nabla f(Z^{\e,\delta}_{T_{i-1}^\e} ) , V_{i} (Z^{\e,\delta}_{T_{i-1}^\e} )\Delta_i T^\e\rangle
= \frac{\e^2}{\delta^2}\sum_{i=2}^{p}\sum_{j=1}^{i-1}    \mu_i\mu_j \langle \nabla g_i(x), V_j(x) \rangle +
 \sum_{i=2}^{p}\sum_{j=1}^{i-1}    \mu_i \E_{x,1}R^{\e,\delta}_{ij}.
\ea
Notice that only indices $j<i$ appear in our analysis:
the value of $V_i$ at the beginning of the $i$-th run is affected by the displacements
generated by the preceding vector fields $V_1,\dots,V_{i-1}$.
This is the point at which the prescribed cyclic ordering enters the limiting drift.


To estimate the remainder, we note that
for all $x,y,z\in\bR^d$
\ba
|\langle\nabla g_i(y+ z), V_j(y) \rangle - \langle \nabla g_i(x), V_j(x)\rangle|
&\leq
|\langle\nabla g_i(y+ z) - \nabla g_i(y), V_j(y) \rangle |\\
&+|\langle\nabla g_i(y), V_j(y) \rangle - \langle \nabla g_i(x), V_j(x) \rangle|\\
&\leq C(|z|  + |x-y| ),
\ea
where the constant $C$ depends only on $\|\nabla^k f\|$, $k=1,2$, and $\|\nabla^k V\|$, $k=0,1$.
Thus, taking into account the Kac scaling $\e/\delta^2\to \varkappa^2$ we get
\ba
\E_{x,1} | R^{\e,\delta}_{ij}|
&\leq_C
\frac{\e}{\delta^2}
\E_{x,1}\Big[\Delta_j T^\e
\int_0^1 \Big(\frac{\theta \Delta_j T^\e}{\delta }|V_j(Z^{\e,\delta}_{T_{j-1}^\e})|
+ |Z^{\e,\delta}_{T_{j-1}^\e}-x|\Big)\,\di \theta \Big]\\
& \leq_C
\frac{\e}{\delta^3 }\E (\Delta_j T^\e)^2
+ \frac{\e}{\delta^2 }\E\Delta_j T^\e \E_{x,1}  |Z^{\e,\delta}_{T_{j-1}^\e}-x| \\
&\leq_C \frac{\e^3 }{\delta^3}
+ \frac{\e^2}{\delta^2 }\sum_{k=1}^{j-1}\E_{x,1}|Z^{\e,\delta}_{T_{k}^\e}-Z^{\e,\delta}_{T_{k-1}^\e}|\\
&\leq_C  \frac{\e^3 }{\delta^3}
+ \frac{\e^3}{\delta^4}
\\
&\leq_C  \e .
\ea
b) We estimate the second term in the last sum in \eqref{e:f}. For each $i=1,\dots,p$ and $x,y,z\in\bR^d$ we have
\ba
\Big|\langle \nabla^2 f(y + z) V_i (y), V_i (y)\rangle  - \langle \nabla^2 f(x) V_i(x), V_i(x)\rangle\Big|
&\leq
\Big|\langle \nabla^2 f(y + z) V_i(y), V_i(y)\rangle - \langle \nabla^2 f(y) V_i(y), V_i (y)\rangle\Big|\\
&+\Big|\langle \nabla^2 f(y) V_i (y), V_i (y)\rangle - \langle \nabla^2 f(x) V_i (x), V_i (x)\rangle\Big|.
\ea
Note that $z\mapsto \langle \nabla^2 f(y + z) V_i(y), V_i(y)\rangle$ is bounded and smooth.
We use here the weaker H\"older-type bound with exponent $\rho\in (0,1)$ in order to estimate
the remainder using only the assumed $(2+\rho)$-moments of the run times.
We have
\ba
\Big|\langle \nabla^2 f(y + z) V_i(y), V_i(y)\rangle - \langle \nabla^2 f(y) V_i(y), V_i (y)\rangle\Big|
\leq_C |z|\wedge 1
\leq_C  |z|^\rho\wedge 1
\leq_C  |z|^{\rho},
\ea
where the constant in this estimate depends only on $\|\nabla^k f\|$, $k=2,3$, and $\|V\|$.
Analogously, since $y\mapsto \langle \nabla^2 f(y) V_i (y), V_i (y)\rangle$ is bounded and smooth, we have
\ba
\Big|\langle \nabla^2 f(y) V_i (y), V_i (y)\rangle - \langle \nabla^2 f(x) V_i (x), V_i (x)\rangle\Big|
\leq_C|x-y|^{\rho},
\ea
where the constant depends on $\|\nabla^k f\|$, $k=2,3$, and $\|\nabla^k V\|$, $k=0,1$.
Therefore,
\ba
& \frac{1}{\delta^2}\sum_{i=1}^p
\E_{x,1} \int_0^1  \Big\langle \nabla^2 f\Big(Z^{\e,\delta}_{T_{i-1}^\e}
+  \frac{\theta}{\delta}V_i (Z^{\e,\delta}_{T_{i-1}^\e} )\Delta_i T^\e\Big) V_i (Z^{\e,\delta}_{T_{i-1}^\e} )\Delta_i T^\e,
V_i (Z^{\e,\delta}_{T_{i-1}^\e} )\Delta_i T^\e\Big\rangle(1-\theta)\,\di \theta\\
&=  \frac{\e^2}{2\delta^2}\sum_{i=1}^p \lambda_i^2
\langle \nabla^2 f(x)V_i (x), V_i (x)\rangle
+\sum_{i=1}^p \E_{x,1} R^{\e,\delta}_i,
\ea
where for each $i=1,\dots,p$
\ba
\E_{x,1} |R^{\e,\delta}_i|
&\leq_C
\frac{1}{\delta^2}
\E_{x,1} \int_0^1 (\Delta_i T^\e)^2\Big(
\Big| \frac{\theta}{\delta}V_i (Z^{\e,\delta}_{T_{i-1}^\e})\Delta_i T^\e\Big|^\rho
+ |Z^{\e,\delta}_{T_{i-1}^\e} - x |^\rho
\Big)(1-\theta)\,\di \theta\\
&\leq_C
\frac{1}{\delta^{2+\rho}}
\E_{x,1} (\Delta_i T^\e)^{2+\rho}
+
\frac{1}{\delta^{2}}
\E_{x,1}\Big[ (\Delta_i T^\e)^{2}
\sum_{j=1}^{i-1}|Z^{\e,\delta}_{T_{j}^\e} - Z^{\e,\delta}_{T_{j-1}^\e}|^\rho\Big]\\
&\leq_C
\frac{1}{\delta^{2+\rho}}
\E_{x,1} (\Delta_i T^\e)^{2+\rho}
+
\frac{1}{\delta^{2+\rho}}\sum_{j=1}^{i-1}
\E_{x,1}\Big[ (\Delta_i T^\e)^{2}(\Delta_j T^\e)^\rho\Big]\\
&\leq_C \frac{\e^{2+\rho}}{\delta^{2+\rho}}
\leq_C  \e^{1+\frac{\rho}{2}}.
\ea
2. Combining the results from a) and b) we get
\ba
\label{e:opA1}
\frac{\E_{x,1}f(Z^{\e,\delta}_{T_p^\e}) -f(x)}{\E_1 T^\e_p}
&=\frac{1}{\mu\e}\frac{\e^2}{\delta^2}
\Big(
\sum_{i=2}^{p}\sum_{j=1}^{i-1}
\mu_i\mu_j \langle \nabla \langle \nabla f(x),  V_i (x)\rangle, V_j(x) \rangle
+\frac{1}{2}  \sum_{i=1}^{p}  \lambda_i^2
\langle \nabla^2 f(x)V_i (x), V_i (x)\rangle
+\E_{x,1} R^{\e,\delta}\Big) \\
&=\frac{\varkappa^2}{\mu}
\Big(
\sum_{j<i} \mu_i\mu_j \langle \nabla \langle \nabla f(x),  V_i (x)\rangle, V_j(x) \rangle
+\frac{1}{2}  \sum_i  \lambda_i^2 \langle \nabla^2 f(x)V_i (x), V_i (x)\rangle \Big)
+ r^{\e,\delta}_{x,1},
\ea
where
\ba
\lim_{\e,\delta\to 0}|r^{\e,\delta}_{x,1}|=0.
\ea
Thus, before using the centering assumption $\mathbf{A}_{V,\mu}$ a second time,
the limiting second-order contribution contains both the individual quadratic
terms associated with each vector field and cross terms generated by successive vector
fields within the cycle.

Recall the identity
\ba
\langle \nabla \langle \nabla f,  V_i\rangle, V_j \rangle
= \langle \nabla^2 f  V_i , V_j \rangle + \langle \nabla f,  D V_i[V_j] \rangle,
\ea
where $D V_i[V_j]=(DV_i)V_j$ is the directional derivative of $V_i$ in direction $V_j$ defined in \eqref{e:DVV}.
We can rewrite
\ba
\label{e:H1}
\frac{\E_{x,1}f(Z^{\e,\delta}_{T_p^\e}) -f(x)}{\E_1 T^\e_p}
&= \frac{\varkappa^2}{\mu} \sum_{j<i} \mu_i\mu_j \langle \nabla^2 f(x)  V_i(x) , V_j(x) \rangle
+ \frac{\varkappa^2}{\mu} \sum_{j<i} \mu_i\mu_j  \langle \nabla f,  D V_i[V_j](x) \rangle\\
&+\frac{\varkappa^2}{2\mu} \sum_{i}  \lambda_i^2 \langle \nabla^2 f(x)V_i(x), V_i(x)\rangle
+r^{\e,\delta}_{x,1}
\ea
The Hessian cross terms can be simplified once more using the centering assumption $\mathbf{A}_{V,\mu}$:
\ba
\label{e:H2}
0=\Big\langle \nabla^2 f \sum_{i=1}^{p} \mu_i V_i,  \sum_{j=1}^{p} \mu_j V_j\Big\rangle
=\sum_{i} \mu_i^2 \langle \nabla^2 f V_i, V_i\rangle
+ 2 \sum_{j<i} \mu_i\mu_j \langle \nabla^2 f V_i, V_j\rangle,
\ea
Substituting \eqref{e:H2} into \eqref{e:H1} replaces the second moments $\lambda_i^2=\sigma_i^2+\mu_i^2$ by the variances $\sigma_i^2$,
leaving the order-dependent first-order term unchanged, so that
\ba
\label{e:AA}
\frac{\E_{x,1}f(Z^{\e,\delta}_{T_p^\e}) -f(x)}{\E_1 T^\e_p}
&=
\frac{\varkappa^2}{2\mu} \sum_i  \sigma_{i}^2
\langle \nabla^2 f(x)V_i(x), V_i(x)\rangle
+\frac{\varkappa^2}{\mu} \sum_{j<i}
\mu_i\mu_j  \langle \nabla f(x),  D V_i[V_j](x) \rangle
+r^{\e,\delta}_{x,1}\\
&= Af(x) + +r^{\e,\delta}_{x,1},
\ea
with the operator $A$ defined in \eqref{e:opA}.

\smallskip
\noindent
3.
It remains to show that the limiting operator is independent of the initial direction $\iota$.
The preceding remainder estimates are uniform in $x$ and $\iota$,
so it suffices to verify that the limiting differential operator is invariant under cyclic shifts of the indices.

Let
\ba
Af=A_1 f:=\frac{\varkappa^2}{2\mu} \sum_i  \sigma_{i}^2
\langle \nabla^2 fV_i , V_i \rangle
+\frac{\varkappa^2}{\mu} \sum_{j<i}
\mu_i\mu_j  \langle \nabla f ,  D V_i[V_j]  \rangle
\ea
be given in \eqref{e:AA}. Let the initial direction be $\iota=2$, and let
 \ba
A_2 f & :=\frac{\varkappa^2}{2\mu} \sum_i  \sigma_{1+i\mathrm{mod}\,p}^2
\langle \nabla^2 f(x)V_{1+i\mathrm{mod}\,p}(x), V_{1+i\mathrm{mod}\,p}(x)\rangle\\
&+\frac{\varkappa^2}{\mu}
\sum_{j<i} \mu_{1+i\mathrm{mod}\,p}\mu_{1+j\mathrm{mod}\,p}  \langle \nabla f,
D V_{1+i\mathrm{mod}\,p}[V_{1+j\mathrm{mod}\,p}] \rangle
\ea
be the corresponding limiting operator with all indices cyclically shifted by 1.
It is clear that the second order parts of $A_1$ and $A_2$ coincide. Therefore,
\ba
A_2f - A_1 f &=\sum_{j=2}^{p} \mu_1\mu_j \Big(\langle \nabla f,  D V_1[V_j] \rangle - \langle \nabla f,  D V_j[V_1] \rangle\Big)\\
&=\mu_1 \Big\langle\nabla f, DV_1\Big[ \sum_{j=2}^{p} \mu_jV_j   \Big]  -  \sum_{j=2}^{p} \mu_j DV_j[V_1]    \Big\rangle.
\ea
Since by assumption $\mathbf{A}_{V,\mu}$
\ba
 \sum_{j=2}^{p} \mu_j V_j=-\mu_1 V_1,
\ea
we get
\ba
  DV_1\Big[ \sum_{j=2}^{p} \mu_jV_j   \Big] =-\mu_1 DV_1[V_1]
\ea
and
\ba
\sum_{j=2}^{p} \mu_j DV_j[V_1] = \Big(D\sum_{j=2}^{p} \mu_jV_j\Big)[V_1]
=-\mu_1 DV_1[V_1].
\ea
Therefore,
\ba
A_2f - A_1f &=\mu_1 \Big\langle\nabla f, -DV_1 [ V_1 ] +   DV_1 [ V_1 ]    \Big\rangle\equiv  0.
\ea
Repeating the same argument for successive cyclic shifts shows that the limiting operator
is independent of the initial direction $\iota$. Notice that this invariance concerns cyclic shifts only.
Example \ref{exa:UV} illustrates that changing the cyclic ordering itself may change the limiting drift.
\end{proof}

\end{document}
